\documentclass[11pt]{article}
\usepackage{preamble}

\newcommand{\IconCheck}{[CHECK]}
\newcommand{\IconMic}{[MIC]}
\newcommand{\IconClipboard}{[CLIPBOARD]}

\newcommand{\phasetag}[1]{%
  {\small\itshape Tag: #1\par}
}

\newcommand{\phaseitems}[1]{%
  \begin{itemize}[leftmargin=1.5em,itemsep=2pt,topsep=2pt]
  #1
  \end{itemize}%
}

\newcommand{\teacheranswerbox}[1]{%
  \begin{callout}{\IconCheck\ ANSWER / ANTICIPATED TALK}{calloutgreen}
  \small #1
  \end{callout}
}

\begin{document}

\packetheading{Lesson 5-4 · Period 1 · Teacher Edition}{Solving Radical Equations --- Foundation}

\sectionbanner{OBJECTIVES}
{\small
\renewcommand{\arraystretch}{1.25}
\noindent\begin{tabular}{|p{1.8in}|p{4.8in}|}
\hline
\textbf{Math Objective} & Solve radical equations by isolating the radical and raising both sides to the appropriate power; identify and reject extraneous solutions. \\ \hline
\textbf{Language Objective} & Explain in writing, using sentence frames, why squaring both sides can introduce extraneous solutions and how to check. \\ \hline
\textbf{Essential Understanding} & Squaring both sides is NOT a reversible operation --- it can create solutions that don't satisfy the original equation. Every candidate must be checked. \\ \hline
\end{tabular}
}

\sectionbanner{TOPIC GOALS / ESSENTIAL QUESTION / MATERIALS}
{\small
\renewcommand{\arraystretch}{1.25}
\noindent\begin{tabular}{|p{1.8in}|p{4.8in}|}
\hline
\textbf{Topic Goals} & Solve single-radical equations; detect extraneous solutions. \\ \hline
\textbf{Essential Question} & Why does squaring both sides sometimes produce solutions that don't work, and how do we catch them? \\ \hline
\textbf{Materials} & Student packet, pencil. Teacher models Example 1, then gradual-release into Example 3 (extraneous). \\ \hline
\end{tabular}
}

\begin{callout}{\IconPin\ PRE-FLIGHT NOTE (read before class)}{calloutpink}
Explore is 6 bank items: Try-It 1 \(\rightarrow\) Try-It 3 \(\rightarrow\) \#22 \(\rightarrow\) \#8 \(\rightarrow\) \#15 \(\rightarrow\) \#4. Paced for \(\sim\)5--6 min per item.\\
If pairs finish early, push them to the Optional Reinforcement: Model \& Discuss Parts B + C (verbal/TPS), then Practice \#21.\\
If pace slows, cut \#4 (extension) --- it's the DOK-2 stretch; the core DOK-1 arc (Try-It 1 \(\rightarrow\) Try-It 3 \(\rightarrow\) \#22 \(\rightarrow\) \#8 \(\rightarrow\) \#15) is the minimum viable set.
\end{callout}

\phasetag{notice \& wonder}
\frameworkphaseheader{Do Now}{2}{5}{%
\phaseitems{
  \item Project the starter (or use printed Do Now sheet).
  \item Circulate silently. Do NOT teach during this phase.
  \item Collect at 5 min sharp.
}%
}{%
\phaseitems{
  \item Write 1 notice + 1 wonder.
  \item Write one sentence about what they would do first to solve a radical equation.
}%
}{%
\phaseitems{
  \item (None during the phase --- teacher harvests answers in Launch.)
}%
}{Solo teacher: circulate silently. Resist the urge to give hints.}

\bankitem{Do Now (5-4-savvas-model-discuss-lesson-5-4-launch)}{%
\textbf{EXPLORE \& REASON}

\textbf{A.} Solve \(3(a + 1)^2 + 2 = 11\). Use at least two different methods.

\textbf{B.} Try each of the methods you used in part (a) to solve \(3\sqrt{a + 1} + 2 = 11\).

\textbf{C.} Generalize Which of the methods is better suited for solving an equation with a radical? What problems arise when using the other method?
}{}

\begin{callout}{\IconCheck\ ANTICIPATED STUDENT RESPONSES}{calloutgreen}
NOTICE (typical): ``There's a square root.'' ``It's not a normal equation.'' ``The variable is under a radical sign.''\\
WONDER (typical): ``How do you undo a square root?'' ``Do you square both sides?'' ``Could there be no solution?''\\
SENTENCES: students might write ``square both sides'' or ``isolate the radical first.''\\
KEY MOVE in Launch: surface that squaring both sides is a one-way operation --- it can create solutions that don't work in the original. That's why CHECK is non-negotiable.
\end{callout}

\phasetag{teacher models Example 1 + Example 3}
\frameworkphaseheader{Launch}{1-2}{12}{%
\phaseitems{
  \item Bridge from Do Now: ``You noticed the square root. To undo it, we square both sides --- but that move can create fake solutions, so we always check.''
  \item Project Example 1 (Savvas Topic 5-4). Think aloud: isolate the radical, square both sides, solve, check the candidate in the ORIGINAL equation.
  \item Project Example 3. Think aloud: isolate, square, solve, check --- one candidate fails the check (extraneous). Cross it out.
  \item Pause 1\(\times\) for 30-sec partner-check: ``In Example 1, what step would catch a fake solution?''
  \item Do NOT solve Try-Its. Release promptly at minute 17.
}%
}{%
\phaseitems{
  \item Watch silently through Example 1.
  \item During partner-check: identify the check step.
  \item Copy the Isolate--Raise--Check Rules card into their page margin.
  \item Note the COMMON ERROR warning: always substitute into the ORIGINAL, not the squared form.
}%
}{%
\phaseitems{
  \item ``What must I do before squaring?''
  \item ``After squaring, how do I verify my answer?''
  \item ``In Example 3, why is one solution extraneous?''
}%
}{Project, think aloud. Do NOT give students Try-Its to work until minute 17.}

\begin{callout}{\IconMic\ TEACHER SCRIPT}{calloutblue}
1. ``From the Do Now, you noticed the radical. The first move is always: isolate the radical on one side.''\\
2. Example 1 walk-through: ``Watch me. I isolate the radical, then square both sides. Now I have a polynomial equation I can solve normally. Then I CHECK my answer in the ORIGINAL --- not the squared form.''\\
3. Example 3 walk-through: ``This time I get two candidates. I check both. One works; one does NOT. That failing one is extraneous --- I cross it out and write `extraneous' next to it.''\\
4. Common-error flag: ``Never skip the check. Squaring is not reversible --- it can create solutions that never existed in the original equation.''\\
5. Release: ``Now you try. Try-It 1 asks you to isolate and solve. Try-It 3 asks you to find and reject an extraneous solution.''
\end{callout}

\phasetag{Think-Pair-Share + Practice}
\frameworkphaseheader{Explore}{1-2}{33}{%
\phaseitems{
  \item 5 laps circulation across 33 min (\(\sim\)6--7 min each).
  \item Lap 1: Try-It 1 --- press pairs to isolate first before squaring.
  \item Lap 2: Try-It 3 --- press pairs to show the check step explicitly.
  \item Lap 3: \#22 + \#8 --- check that pairs are substituting into the ORIGINAL equation.
  \item Lap 4: \#15 --- verify pairs identify and label the extraneous solution.
  \item Lap 5: \#4 (extension) --- stretch. Cut if pace slows.
  \item Do not give answers. Use sentence frame to press justification.
}%
}{%
\phaseitems{
  \item Pairs move in order through 6 items: 2 Try-Its + 4 Practice.
  \item For each item: isolate the radical FIRST, then raise, then check.
  \item For Practice \#15 (identify extraneous): show the check step and cross out the extraneous solution.
  \item Justify with sentence frame before writing.
}%
}{%
\phaseitems{
  \item Did you isolate the radical before raising both sides?
  \item Which equation did you substitute into to check --- original or squared?
  \item If a candidate doesn't satisfy the original, what do you write?
  \item Can you use the sentence frame to explain your answer?
  \item For \#4: show all steps including the check.
}%
}{Solo teacher: 5 laps. Press pairs to justify using the rule card. No answers.}

\bankitem{Try It 1}{%
Solve each radical equation. Check your solution by substituting into the given equation or graphing.

\textbf{a.} \(\sqrt{x - 2} + 3 = 5\)

\textbf{b.} \(\sqrt[3]{x - 1} = 2\)
}{}

\teacheranswerbox{Answer key (Savvas TE): A. \(x=6\); B. \(x=9\).}

\bankitem{Try It 3 --- extraneous solution}{%
Solve each radical equation. Check for any extraneous solutions.

\textbf{a.} \(x = \sqrt{7x + 8}\)

\textbf{b.} \(x + 2 = \sqrt{x + 2}\)
}{}

\teacheranswerbox{Answer key (Savvas TE): A. \(x=8\); \(x=-1\) is an extraneous solution. B. \(x=-1\), \(x=-2\).}

\bankitem{Practice \#22 --- Solve radical equation}{%
Solve each radical equation. \(\sqrt{4x}=11\)
}{}

\teacheranswerbox{Answer key (from source): \(x=30.25\).}

\bankitem{Practice \#8 --- Solve radical equation}{%
\(\sqrt{8x + 9} = x\)
}{}

\teacheranswerbox{Savvas TE answer: \(x=-1\). Teacher note: squaring gives candidates \(x=9\) and \(x=-1\); \(x=-1\) is extraneous, so the valid solution to the original equation is \(x=9\).}

\bankitem{Practice \#15 --- Identify extraneous}{%
Generalize Explain how to identify an extraneous solution for an equation containing a radical expression.
}{}

\teacheranswerbox{Substitute each potential solution into the original equation. If it makes the equation true, it is a real solution.}

\bankitem{Practice \#4 --- Solve and verify}{%
Error Analysis Nazim said that \(-3\) and \(6\) are the solutions to \(\sqrt{3x + 18} = x\). What error did Nazim make?
}{}

\teacheranswerbox{Nazim didn't check for extraneous solutions. The real solution is \(6\); \((-3)\) is extraneous.}

\phasetag{academic conversation}
\frameworkphaseheader{Share / Summary}{2}{5}{%
\phaseitems{
  \item Cold-call 2 pairs to read their sentence-frame sentence.
  \item Write one student-generated sentence on board.
  \item Preview P2: ``Next class we'll tackle equations with two radicals or a radical and a binomial --- the same Isolate--Raise--Check logic applies.''
}%
}{%
\phaseitems{
  \item Presenting pair reads their sentence.
  \item Class signals thumbs up/sideways.
  \item Note: next class extends this to more complex radical equations.
}%
}{%
\phaseitems{
  \item What's a sentence that captures today's rule about extraneous solutions?
  \item Why can't we trust a candidate solution without checking it?
}%
}{Cold-call 2 pairs. Keep it brief --- 5 min.}

\begin{sentenceframebox}
Pair share: ``I isolated the radical to get \_\_\_. Squaring both sides gives \_\_\_. My candidate solutions are \_\_\_. When I check in the ORIGINAL equation, \_\_\_ is valid and \_\_\_ is extraneous because \_\_\_.''\\
Whole class: one pair reads their sentence. Class signals thumbs up/sideways.
\end{sentenceframebox}

\phasetag{summary recap (not graded)}
\frameworkphaseheader{Exit Ticket}{1-2}{2}{%
\phaseitems{
  \item Collect at door.
  \item Scan the fill-in items --- misconceptions about the check step or extraneous solutions tell you what to re-hit in P2 Do Now.
}%
}{%
\phaseitems{
  \item Complete the fill-in stems.
  \item Be honest about confusion.
}%
}{%
\phaseitems{
  \item Did students correctly identify what an extraneous solution is and why it arises?
}%
}{Collect. No grade.}

\begin{callout}{\IconClipboard\ EXIT TICKET --- Student Form}{calloutyellow}
To solve a radical equation I first \_\_\_.\\
I raise both sides to power \_\_\_ because \_\_\_.\\
An extraneous solution happens when \_\_\_.\\
One biggest thing I learned today: \rule{2.6in}{0.4pt}
\end{callout}

\sectionbanner{OPTIONAL REINFORCEMENT (not collected)}
\textit{If you finish Explore early or want extra practice before Period 2.}

\begin{callout}{\IconSearch\ OPTIONAL · Model \& Discuss}{calloutblue}
Return to the Do Now equation. With your partner:\\
\\
B) Use Structure. Why can squaring both sides introduce a solution that doesn't work in the original equation?\\
\\
C) Examine at least two equations you solved in Explore. For each, explain in your own words why checking is a required step and not optional.
\end{callout}

\bankitem{Optional \#1 (Savvas Practice)}{%
Solve each radical equation. \(\sqrt[3]{x}+8=13\)
}{}

\teacheranswerbox{Answer key (from source): \(x=125\).}

\clearpage

\begin{callout}{\IconClock\ PERIOD A \& F --- 65-MIN CLASS NOTES}{calloutpink}
Both sections get 65 min for P1. Use the extra 10 min on Explore, NOT Launch.\\
Suggested add: project Savvas Practice \#15 and \#4 on screen for 2 more TPS rounds if pairs finish fast.\\
Or: extend Model \& Discuss Parts B + C with ``write one more radical equation that has an extraneous solution --- then solve and reject it.''
\end{callout}

\begin{callout}{\IconPuzzle\ MODIFICATIONS / SUPPORT FOR IEP STUDENTS}{calloutiep}
\textbullet\ Provide the Isolate--Raise--Check Rules card as a sticker/cut-out on their desk.\\
\textbullet\ For Try-It 1, allow students to use a step-by-step graphic organizer (isolate \(\rightarrow\) raise \(\rightarrow\) check).\\
\textbullet\ Accept verbal sentence-frame completions in lieu of written.
\end{callout}

\begin{callout}{\IconGlobe\ SUPPORTS FOR ELL STUDENTS}{calloutell}
\textbullet\ Sentence frames printed on student packet (don't skip).\\
\textbullet\ Word card: ``extraneous solution'' = a candidate answer that does not satisfy the original equation; must be rejected.\\
\textbullet\ ``isolate'' = get the radical alone on one side of the equation.\\
\textbullet\ ``raise to the index power'' = square both sides for a square root; cube both sides for a cube root.\\
\textbullet\ Pair ELL students with English-dominant partners for TPS.
\end{callout}

\end{document}
